% !TeX root = ../main.tex
\chapter{Conclusion}
\label{chap:conclusion}

Ce laboratoire a permis de caractériser et de comparer expérimentalement le comportement statique et dynamique de deux types de capteurs de température courants : la thermistance CTN et la sonde RTD Pt100.

\section{Bilan des objectifs}
\label{sec:conclusion_objectifs}

Tous les objectifs fixés ont été atteints :
\begin{itemize}[leftmargin=1.2cm]
    \item \textbf{Mesure ambiante CTN :} $T_{\text{amb}} = \qty{21.54 \pm 0.80}{\degreeCelsius}$ déduite par interpolation locale de la datasheet.
    \item \textbf{Auto-échauffement :} Le pont sous \qty{14}{\volt} avec $R = \qty{462.35}{\kilo\ohm}$ limite l'échauffement Joule à $\Delta T \approx \qty{0.09}{\degreeCelsius}$, négligeable devant la tolérance de $\pm \qty{0.8}{\degreeCelsius}$.
    \item \textbf{Acquisition LabVIEW et filtrage :} Linéarisation logicielle par table de correspondance. L'essai au sèche-cheveux montre que le suréchantillonnage lisse le signal numérique, mais que le bruit résiduel ($\sigma \approx \qty{3}{\degreeCelsius}$) est dominé par les turbulences réelles du fluide et la tenue manuelle.
    \item \textbf{Étalonnage de la Pt100 :} Étalonnage au puits sec donnant $R_0 = \qty{100.10 \pm 0.03}{\ohm}$ à \qty{0.00}{\degreeCelsius}.
    \item \textbf{Validation du 4 fils Kelvin :} Mise en évidence du biais en 2 fils ($2 R_L = \qty{0.09}{\ohm} \implies \Delta T = \qty{+0.23}{\degreeCelsius}$), démontrant la nécessité du montage 4 fils pour une sonde de faible résistance nominale.
    \item \textbf{Dynamique thermique en eau chaude :} Modélisation au premier ordre ($\text{RMSE} = \qty{0.09}{\degreeCelsius}$) avec $\tau = \qty{5.80 \pm 0.05}{\second}$ et $t_{95\%} = \qty{17.37 \pm 0.15}{\second}$, gouvernés par l'inertie thermique de la gaine inox.
\end{itemize}

\section{Comparaison technologique et bilan}
\label{sec:conclusion_synthese}

Ce laboratoire illustre les compromis métrologiques classiques entre les deux technologies :
\begin{itemize}[leftmargin=1.2cm]
    \item \textbf{Thermistance CTN :} Très forte sensibilité ($\approx \qty{4.2}{\kilo\ohm\per\degreeCelsius}$), facilitant son conditionnement. En contrepartie, elle est très non-linéaire et présente une tolérance plus large ($\pm \qty{0.8}{\degreeCelsius}$ sans étalonnage individuel).
    \item \textbf{Sonde RTD Pt100 :} Sensibilité plus faible ($\qty{0.385}{\ohm\per\degreeCelsius}$) imposant un câblage 4 fils Kelvin, mais excellente linéarité et haute précision (classe A : $\pm \qty{0.15}{\degreeCelsius}$), idéale comme capteur étalon.
\end{itemize}
Les deux capteurs concordent parfaitement à température ambiante ($21.54\,^\circ\text{C}$ vs $22.16\,^\circ\text{C}$), résultat validé par le critère de compatibilité ($C = 0.76 \le 2$).
